% --- Archon physics formalization source begin ---
% archon:physics
% archon:covers PhyXMiniProblems/problem_phyx_mini_0072.lean
% archon:source-report reports/phyx_mini/problem_phyx_mini_0072.source.json

\chapter{Physics problem phyx\_mini\_0072}
\label{ch:PhyXMiniProblems_problem_phyx_mini_0072}

\paragraph{Problem source.}
\#\# Physical scenario

Shown from above in figure is one corner of a rectangular box filled with water. A laser beam starts 10\textbackslash{} \textbackslash{}mathrm\{cm\} from side A of the container and enters the water at position x. You can ignore the thin walls of the container.

\#\# Question

Find the minimum value of x for which the laser beam passes through side B and emerges into the air.

\#\# Answer choices

- A: A: 20cm

- B: B: 15.0cm

- C: C: 18.2cm

- D: D: 15.3cm

\#\# Figure caption (auxiliary; use the image as primary evidence)

The image depicts a diagram with the following components:

1. A rectangular area labeled as "Water (top view)" in the center. This area is shaded light blue and represents water.

2. The water has two labeled sides: "Side A," along the bottom, and "Side B," along the right edge.

3. To the left of the water, there is a 3D object resembling a cuboid with a circle on its face, positioned at an angle and connected to the water's edge with a purple arrow.

4. The purple arrow extends from the circle on the cuboid to the left edge of Side A, and it is labeled with "10 cm."

5. The diagram includes the variable \textbackslash{}( x \textbackslash{}), indicated by a vertical arrow along the left side of the water area, suggesting a measurement related to the water.

No additional text or annotations are present in the diagram.

\#\# Dataset metadata

Geometrical Optics; Physical Model Grounding Reasoning, Spatial Relation Reasoning

\paragraph{Recorded answer/context.}
C: C: 18.2cm

\paragraph{Figure/image path.}
/root/proposal\_for\_physic/science-mango/phyx\_mini\_run/phyx\_data/test\_image/72.png

\paragraph{Formalization target.}
create a compiling Lean file with sorry bodies at `PhyXMiniProblems/problem\_phyx\_mini\_0072.lean`. The Lean declarations must preserve the physical quantities, dimensions or dimensional roles, figure labels, governing-law hypotheses, and final relation expressed by this problem.
Use Mathlib/Physlib names found through LeanExplore where available. If a physics API is missing, introduce faithful local abstractions rather than scalar placeholder aliases.

\begin{theorem}[Physics formalization target]
\label{thm:physics:phyx_mini_0072:target}
\lean{PhyXMiniProblems.ProblemPhyXMini0072.problem_phyx_mini_0072}
\uses{def:physics:phyx-mini-0072:phyxminiproblems-problemphyxmini0072-rectangularwaterboxsetup, def:physics:phyx-mini-0072:phyxminiproblems-problemphyxmini0072-hasdepictedrectangularlayout, def:physics:phyx-mini-0072:phyxminiproblems-problemphyxmini0072-matchesproblemandfigurereadouts, def:physics:phyx-mini-0072:phyxminiproblems-problemphyxmini0072-hasphysicalopticalparameters, def:physics:phyx-mini-0072:phyxminiproblems-problemphyxmini0072-satisfieswateraircriticalanglelaw, def:physics:phyx-mini-0072:phyxminiproblems-problemphyxmini0072-admissibleentrydisplacementsincentimeters, def:physics:phyx-mini-0072:phyxminiproblems-problemphyxmini0072-criticalthresholdentrydisplacementincentimeters, def:physics:phyx-mini-0072:phyxminiproblems-problemphyxmini0072-matchesentrydisplacementchoice}
The assigned autoformalize agent should translate this physics problem into Lean declarations in the covered file, with theorem and lemma proof bodies written as `by sorry`.
\end{theorem}
\begin{proof}
This is an autoformalization task, not a proof task. Produce statements that can later be proved by the physics prover without weakening the physical model.
\end{proof}

% --- Archon declaration topology begin ---
\section{Declaration topology}
\begin{definition}[Length Quantity]
  \label{def:physics:phyx-mini-0072:phyxminiproblems-problemphyxmini0072-lengthquantity}
  \lean{PhyXMiniProblems.ProblemPhyXMini0072.LengthQuantity}
  A nonnegative physical length, represented independently of unit choice
\end{definition}
\begin{proof}
  This is the declared model definition of Length Quantity.
\end{proof}
\begin{definition}[length Readout]
  \label{def:physics:phyx-mini-0072:phyxminiproblems-problemphyxmini0072-lengthreadout}
  \lean{PhyXMiniProblems.ProblemPhyXMini0072.lengthReadout}
  \uses{def:physics:phyx-mini-0072:phyxminiproblems-problemphyxmini0072-lengthquantity}
  The scalar readout of a physical length in a selected length unit
\end{definition}
\begin{proof}
  This is the declared model definition of length Readout.
\end{proof}
\begin{definition}[length In Centimeters]
  \label{def:physics:phyx-mini-0072:phyxminiproblems-problemphyxmini0072-lengthincentimeters}
  \lean{PhyXMiniProblems.ProblemPhyXMini0072.lengthInCentimeters}
  \uses{def:physics:phyx-mini-0072:phyxminiproblems-problemphyxmini0072-lengthquantity, def:physics:phyx-mini-0072:phyxminiproblems-problemphyxmini0072-lengthreadout}
  The scalar centimeter readout of a physical length
\end{definition}
\begin{proof}
  This is the declared model definition of length In Centimeters.
\end{proof}
\begin{definition}[Container Side]
  \label{def:physics:phyx-mini-0072:phyxminiproblems-problemphyxmini0072-containerside}
  \lean{PhyXMiniProblems.ProblemPhyXMini0072.ContainerSide}
  The two perpendicular container sides named in the figure
\end{definition}
\begin{proof}
  This is the declared model definition of Container Side.
\end{proof}
\begin{definition}[Boundary Orientation]
  \label{def:physics:phyx-mini-0072:phyxminiproblems-problemphyxmini0072-boundaryorientation}
  \lean{PhyXMiniProblems.ProblemPhyXMini0072.BoundaryOrientation}
  Top-view orientations of the two boundary segments
\end{definition}
\begin{proof}
  This is the declared model definition of Boundary Orientation.
\end{proof}
\begin{definition}[Optical Medium]
  \label{def:physics:phyx-mini-0072:phyxminiproblems-problemphyxmini0072-opticalmedium}
  \lean{PhyXMiniProblems.ProblemPhyXMini0072.OpticalMedium}
  The homogeneous media crossed by the laser ray
\end{definition}
\begin{proof}
  This is the declared model definition of Optical Medium.
\end{proof}
\begin{definition}[Rectangular Water Box Setup]
  \label{def:physics:phyx-mini-0072:phyxminiproblems-problemphyxmini0072-rectangularwaterboxsetup}
  \lean{PhyXMiniProblems.ProblemPhyXMini0072.RectangularWaterBoxSetup}
  \uses{def:physics:phyx-mini-0072:phyxminiproblems-problemphyxmini0072-lengthquantity, def:physics:phyx-mini-0072:phyxminiproblems-problemphyxmini0072-containerside, def:physics:phyx-mini-0072:phyxminiproblems-problemphyxmini0072-boundaryorientation, def:physics:phyx-mini-0072:phyxminiproblems-problemphyxmini0072-opticalmedium}
  The fixed apparatus and material data. The critical angle is measured in water from the normal to a water--air interface
\end{definition}
\begin{proof}
  This is the declared model definition of Rectangular Water Box Setup.
\end{proof}
\begin{definition}[Laser Ray Path]
  \label{def:physics:phyx-mini-0072:phyxminiproblems-problemphyxmini0072-laserraypath}
  \lean{PhyXMiniProblems.ProblemPhyXMini0072.LaserRayPath}
  \uses{def:physics:phyx-mini-0072:phyxminiproblems-problemphyxmini0072-containerside}
  One candidate ray through the two named sides. The entry displacement is the centimeter readout labeled x; all angles are in radians and are measured from the normal of the boundary at which they occur
\end{definition}
\begin{proof}
  This is the declared model definition of Laser Ray Path.
\end{proof}
\begin{definition}[Has Depicted Rectangular Layout]
  \label{def:physics:phyx-mini-0072:phyxminiproblems-problemphyxmini0072-hasdepictedrectangularlayout}
  \lean{PhyXMiniProblems.ProblemPhyXMini0072.HasDepictedRectangularLayout}
  \uses{def:physics:phyx-mini-0072:phyxminiproblems-problemphyxmini0072-lengthincentimeters, def:physics:phyx-mini-0072:phyxminiproblems-problemphyxmini0072-rectangularwaterboxsetup}
  Side labels, orientations, media, and the negligible-wall idealization
\end{definition}
\begin{proof}
  This is the declared model definition of Has Depicted Rectangular Layout.
\end{proof}
\begin{definition}[Matches Problem And Figure Readouts]
  \label{def:physics:phyx-mini-0072:phyxminiproblems-problemphyxmini0072-matchesproblemandfigurereadouts}
  \lean{PhyXMiniProblems.ProblemPhyXMini0072.MatchesProblemAndFigureReadouts}
  \uses{def:physics:phyx-mini-0072:phyxminiproblems-problemphyxmini0072-lengthincentimeters, def:physics:phyx-mini-0072:phyxminiproblems-problemphyxmini0072-rectangularwaterboxsetup}
  The given/calibrated numerical readouts: source distance 10 cm, air index 1.00, and water index 1.33. No value for x occurs in this predicate
\end{definition}
\begin{proof}
  This is the declared model definition of Matches Problem And Figure Readouts.
\end{proof}
\begin{definition}[Has Physical Optical Parameters]
  \label{def:physics:phyx-mini-0072:phyxminiproblems-problemphyxmini0072-hasphysicalopticalparameters}
  \lean{PhyXMiniProblems.ProblemPhyXMini0072.HasPhysicalOpticalParameters}
  \uses{def:physics:phyx-mini-0072:phyxminiproblems-problemphyxmini0072-lengthincentimeters, def:physics:phyx-mini-0072:phyxminiproblems-problemphyxmini0072-rectangularwaterboxsetup}
  Positive optical data and the acute physical branch for the critical ray
\end{definition}
\begin{proof}
  This is the declared model definition of Has Physical Optical Parameters.
\end{proof}
\begin{definition}[Satisfies Water Air Critical Angle Law]
  \label{def:physics:phyx-mini-0072:phyxminiproblems-problemphyxmini0072-satisfieswateraircriticalanglelaw}
  \lean{PhyXMiniProblems.ProblemPhyXMini0072.SatisfiesWaterAirCriticalAngleLaw}
  \uses{def:physics:phyx-mini-0072:phyxminiproblems-problemphyxmini0072-rectangularwaterboxsetup}
  Snell's law for the limiting water--air ray. At the critical angle the transmitted air ray is tangent to side B, so its sine is one
\end{definition}
\begin{proof}
  This is the declared model definition of Satisfies Water Air Critical Angle Law.
\end{proof}
\begin{definition}[Follows Depicted Side Route]
  \label{def:physics:phyx-mini-0072:phyxminiproblems-problemphyxmini0072-followsdepictedsideroute}
  \lean{PhyXMiniProblems.ProblemPhyXMini0072.FollowsDepictedSideRoute}
  \uses{def:physics:phyx-mini-0072:phyxminiproblems-problemphyxmini0072-laserraypath}
  The ray enters through side A and reaches side B, as depicted
\end{definition}
\begin{proof}
  This is the declared model definition of Follows Depicted Side Route.
\end{proof}
\begin{definition}[Has Physical Ray Parameters]
  \label{def:physics:phyx-mini-0072:phyxminiproblems-problemphyxmini0072-hasphysicalrayparameters}
  \lean{PhyXMiniProblems.ProblemPhyXMini0072.HasPhysicalRayParameters}
  \uses{def:physics:phyx-mini-0072:phyxminiproblems-problemphyxmini0072-laserraypath}
  The displacement is nonnegative and all normal-relative ray angles lie on the physical acute branch. The strict upper bounds at side A exclude Lean's totalized tan (π / 2) = 0 endpoint from the geometric model
\end{definition}
\begin{proof}
  This is the declared model definition of Has Physical Ray Parameters.
\end{proof}
\begin{definition}[Satisfies Source To Entry Geometry]
  \label{def:physics:phyx-mini-0072:phyxminiproblems-problemphyxmini0072-satisfiessourcetoentrygeometry}
  \lean{PhyXMiniProblems.ProblemPhyXMini0072.SatisfiesSourceToEntryGeometry}
  \uses{def:physics:phyx-mini-0072:phyxminiproblems-problemphyxmini0072-lengthincentimeters, def:physics:phyx-mini-0072:phyxminiproblems-problemphyxmini0072-rectangularwaterboxsetup, def:physics:phyx-mini-0072:phyxminiproblems-problemphyxmini0072-laserraypath}
  Right-triangle geometry outside side A: x = d tan θ\_air, where d is the normal source distance and both lengths are read in centimeters
\end{definition}
\begin{proof}
  This is the declared model definition of Satisfies Source To Entry Geometry.
\end{proof}
\begin{definition}[Satisfies Perpendicular Side Geometry]
  \label{def:physics:phyx-mini-0072:phyxminiproblems-problemphyxmini0072-satisfiesperpendicularsidegeometry}
  \lean{PhyXMiniProblems.ProblemPhyXMini0072.SatisfiesPerpendicularSideGeometry}
  \uses{def:physics:phyx-mini-0072:phyxminiproblems-problemphyxmini0072-laserraypath}
  Because sides A and B are perpendicular, the water-ray angle from the side-A normal complements its incidence angle from the side-B normal
\end{definition}
\begin{proof}
  This is the declared model definition of Satisfies Perpendicular Side Geometry.
\end{proof}
\begin{definition}[Satisfies Snell Law At Side A]
  \label{def:physics:phyx-mini-0072:phyxminiproblems-problemphyxmini0072-satisfiessnelllawatsidea}
  \lean{PhyXMiniProblems.ProblemPhyXMini0072.SatisfiesSnellLawAtSideA}
  \uses{def:physics:phyx-mini-0072:phyxminiproblems-problemphyxmini0072-rectangularwaterboxsetup, def:physics:phyx-mini-0072:phyxminiproblems-problemphyxmini0072-laserraypath}
  Snell's law at side A for transmission from ambient air into water
\end{definition}
\begin{proof}
  This is the declared model definition of Satisfies Snell Law At Side A.
\end{proof}
\begin{definition}[Emerges Into Air Through Side B]
  \label{def:physics:phyx-mini-0072:phyxminiproblems-problemphyxmini0072-emergesintoairthroughsideb}
  \lean{PhyXMiniProblems.ProblemPhyXMini0072.EmergesIntoAirThroughSideB}
  \uses{def:physics:phyx-mini-0072:phyxminiproblems-problemphyxmini0072-rectangularwaterboxsetup, def:physics:phyx-mini-0072:phyxminiproblems-problemphyxmini0072-laserraypath}
  A water ray emerges through side B on the modeled branch exactly when its incidence angle is no larger than the water--air critical angle
\end{definition}
\begin{proof}
  This is the declared model definition of Emerges Into Air Through Side B.
\end{proof}
\begin{definition}[Admissible Entry Displacement In Centimeters]
  \label{def:physics:phyx-mini-0072:phyxminiproblems-problemphyxmini0072-admissibleentrydisplacementincentimeters}
  \lean{PhyXMiniProblems.ProblemPhyXMini0072.AdmissibleEntryDisplacementInCentimeters}
  \uses{def:physics:phyx-mini-0072:phyxminiproblems-problemphyxmini0072-rectangularwaterboxsetup, def:physics:phyx-mini-0072:phyxminiproblems-problemphyxmini0072-laserraypath, def:physics:phyx-mini-0072:phyxminiproblems-problemphyxmini0072-followsdepictedsideroute, def:physics:phyx-mini-0072:phyxminiproblems-problemphyxmini0072-hasphysicalrayparameters, def:physics:phyx-mini-0072:phyxminiproblems-problemphyxmini0072-satisfiessourcetoentrygeometry, def:physics:phyx-mini-0072:phyxminiproblems-problemphyxmini0072-satisfiesperpendicularsidegeometry, def:physics:phyx-mini-0072:phyxminiproblems-problemphyxmini0072-satisfiessnelllawatsidea, def:physics:phyx-mini-0072:phyxminiproblems-problemphyxmini0072-emergesintoairthroughsideb}
  A displacement is admissible when a physical ray with that displayed readout follows the depicted route and satisfies the geometry, both refraction laws, and the side-B escape criterion. This predicate contains no answer choice and does not assert that any particular displacement is minimal
\end{definition}
\begin{proof}
  This is the declared model definition of Admissible Entry Displacement In Centimeters.
\end{proof}
\begin{definition}[admissible Entry Displacements In Centimeters]
  \label{def:physics:phyx-mini-0072:phyxminiproblems-problemphyxmini0072-admissibleentrydisplacementsincentimeters}
  \lean{PhyXMiniProblems.ProblemPhyXMini0072.admissibleEntryDisplacementsInCentimeters}
  \uses{def:physics:phyx-mini-0072:phyxminiproblems-problemphyxmini0072-rectangularwaterboxsetup, def:physics:phyx-mini-0072:phyxminiproblems-problemphyxmini0072-admissibleentrydisplacementincentimeters}
  The set of centimeter readouts for which the ray can escape at side B
\end{definition}
\begin{proof}
  This is the declared model definition of admissible Entry Displacements In Centimeters.
\end{proof}
\begin{definition}[critical Threshold Entry Displacement In Centimeters]
  \label{def:physics:phyx-mini-0072:phyxminiproblems-problemphyxmini0072-criticalthresholdentrydisplacementincentimeters}
  \lean{PhyXMiniProblems.ProblemPhyXMini0072.criticalThresholdEntryDisplacementInCentimeters}
  \uses{def:physics:phyx-mini-0072:phyxminiproblems-problemphyxmini0072-lengthincentimeters, def:physics:phyx-mini-0072:phyxminiproblems-problemphyxmini0072-rectangularwaterboxsetup}
  The critical-ray candidate threshold. At threshold the water-side angle from the side-A normal is complementary to the side-B critical angle. Snell's law then gives sin θ\_air = (n\_water / n\_air) cos θc, followed by x = d tan θ\_air. Naming this candidate does not assert its minimality
\end{definition}
\begin{proof}
  This is the declared model definition of critical Threshold Entry Displacement In Centimeters.
\end{proof}
\begin{definition}[Answer Choice]
  \label{def:physics:phyx-mini-0072:phyxminiproblems-problemphyxmini0072-answerchoice}
  \lean{PhyXMiniProblems.ProblemPhyXMini0072.AnswerChoice}
  Labels of the four displayed answer choices
\end{definition}
\begin{proof}
  This is the declared model definition of Answer Choice.
\end{proof}
\begin{definition}[answer Entry Displacement In Centimeters]
  \label{def:physics:phyx-mini-0072:phyxminiproblems-problemphyxmini0072-answerentrydisplacementincentimeters}
  \lean{PhyXMiniProblems.ProblemPhyXMini0072.answerEntryDisplacementInCentimeters}
  \uses{def:physics:phyx-mini-0072:phyxminiproblems-problemphyxmini0072-answerchoice}
  The centimeter readout printed beside each answer label
\end{definition}
\begin{proof}
  This is the declared model definition of answer Entry Displacement In Centimeters.
\end{proof}
\begin{definition}[Matches Entry Displacement Choice]
  \label{def:physics:phyx-mini-0072:phyxminiproblems-problemphyxmini0072-matchesentrydisplacementchoice}
  \lean{PhyXMiniProblems.ProblemPhyXMini0072.MatchesEntryDisplacementChoice}
  \uses{def:physics:phyx-mini-0072:phyxminiproblems-problemphyxmini0072-answerchoice, def:physics:phyx-mini-0072:phyxminiproblems-problemphyxmini0072-answerentrydisplacementincentimeters}
  Agreement with an answer displayed to the nearest tenth of a centimeter
\end{definition}
\begin{proof}
  This is the declared model definition of Matches Entry Displacement Choice.
\end{proof}
\begin{lemma}[critical Threshold is Least Admissible]
  \label{lem:physics:phyx-mini-0072:phyxminiproblems-problemphyxmini0072-criticalthreshold-isleastadmissible}
  \lean{PhyXMiniProblems.ProblemPhyXMini0072.criticalThreshold_isLeastAdmissible}
  \uses{def:physics:phyx-mini-0072:phyxminiproblems-problemphyxmini0072-rectangularwaterboxsetup, def:physics:phyx-mini-0072:phyxminiproblems-problemphyxmini0072-hasdepictedrectangularlayout, def:physics:phyx-mini-0072:phyxminiproblems-problemphyxmini0072-matchesproblemandfigurereadouts, def:physics:phyx-mini-0072:phyxminiproblems-problemphyxmini0072-hasphysicalopticalparameters, def:physics:phyx-mini-0072:phyxminiproblems-problemphyxmini0072-satisfieswateraircriticalanglelaw, def:physics:phyx-mini-0072:phyxminiproblems-problemphyxmini0072-admissibleentrydisplacementsincentimeters, def:physics:phyx-mini-0072:phyxminiproblems-problemphyxmini0072-criticalthresholdentrydisplacementincentimeters}
  The ray at the critical threshold is the least admissible entry displacement. This is a derived conclusion from the side-A Snell law, perpendicular-side geometry, side-B critical-angle law, and the acute physical branches
\end{lemma}
\begin{proof}
  The conclusion follows from the governing relations and figure readouts named in the statement of critical Threshold is Least Admissible.
\end{proof}
\begin{lemma}[critical Threshold matches choice C]
  \label{lem:physics:phyx-mini-0072:phyxminiproblems-problemphyxmini0072-criticalthreshold-matches-choice-c}
  \lean{PhyXMiniProblems.ProblemPhyXMini0072.criticalThreshold_matches_choice_C}
  \uses{def:physics:phyx-mini-0072:phyxminiproblems-problemphyxmini0072-rectangularwaterboxsetup, def:physics:phyx-mini-0072:phyxminiproblems-problemphyxmini0072-matchesproblemandfigurereadouts, def:physics:phyx-mini-0072:phyxminiproblems-problemphyxmini0072-hasphysicalopticalparameters, def:physics:phyx-mini-0072:phyxminiproblems-problemphyxmini0072-satisfieswateraircriticalanglelaw, def:physics:phyx-mini-0072:phyxminiproblems-problemphyxmini0072-criticalthresholdentrydisplacementincentimeters, def:physics:phyx-mini-0072:phyxminiproblems-problemphyxmini0072-matchesentrydisplacementchoice}
  For the 10 cm, n\_air = 1.00, and n\_water = 1.33 readouts, the critical threshold is within 0.05 cm of 18.2 cm
\end{lemma}
\begin{proof}
  The conclusion follows from the governing relations and figure readouts named in the statement of critical Threshold matches choice C.
\end{proof}
% --- Archon declaration topology end ---
% --- Archon physics formalization source end ---
